\chapter{Signal Descriptions}
\label{ch:signal-descriptions}

\section{Signal Conventions}

Each signal lists its own assertion convention. Address, data, and bus-type value pins use positive logic, while bus
control, handshake, interrupt, halt, trace, and reset pins generally use active-low conventions. The descriptions below
use "asserted" and "deasserted" to refer to the logical state of a signal according to its signal type.

\section{Address Bus Bits (A0--A23)}

\begin{itemize}
	\item \textbf{Direction:} Output
	\item \textbf{Signal Type:} Active High
\end{itemize}

Connected to the address bus.

These 24 pins declare the address on the bus that the CPU reads from or writes to when \activelow{BAS} (bus access
strobe) is asserted.

\section{Power Pins (VCC1--VCC2, GND1--GND2)}

These pins provide the 5~V power that the CPU requires to function.

The VCC pins must be connected to 5~V $\pm$ 5\%, and GND pins must be connected to a 0~V ground rail.

% TODO: Should I recommend some capacitors etc. to be put on the rail or something like that? Maybe should see what other CPU data sheets say
% category=Manual

\section{Data Bus Bits (D0--D15)}

\begin{itemize}
	\item \textbf{Direction:} Bidirectional
	\item \textbf{Signal Type:} Active High
\end{itemize}

Connected to the data bus.

These 16 pins are driven to a 16-bit value when the CPU is outputting to the bus, and are an input for a 16-bit value
when the CPU is reading from the bus.

\section{Clock Input (CLKI)}

\begin{itemize}
	\item \textbf{Direction:} Input
	\item \textbf{Signal Type:} Rising Edge
\end{itemize}

Progresses the CPU to the next cycle when the input transitions from low to high (rising edge).

Each rising edge of CLKI advances the CPU by one execution phase, except while a bus wait state holds the phase, while
the CPU is halted or waiting (\mnemonic{WAIT}), or while it is held in reset. With no wait states, an ordinary
instruction takes six cycles. Coprocessor calls take longer; see the Instruction Timing appendix of the \archmanual{}.

\section{Instruction Sync Output (SYNC)}

\begin{itemize}
	\item \textbf{Direction:} Output
	\item \textbf{Signal Type:} Active High
\end{itemize}

Asserted only at phase 0 of an instruction fetch or of an exception-unit dispatch slot. It remains asserted if a wait
state holds that phase. It is deasserted in every other phase, in every other coprocessor slot (DMA and maths), in
every DMA continuation slot, and while the CPU is halted, waiting, or in reset.

This pin is useful for debugging.

\section{Bus Access Strobe Output (BAS)}

\begin{itemize}
	\item \textbf{Direction:} Output
	\item \textbf{Signal Type:} Active Low
\end{itemize}

When the CPU is ready to perform a bus operation, \activelow{BAS} is asserted after the address pins, BRW, BAT, PROT,
and any CPU-driven data pins are valid.

\activelow{BAS} remains asserted until the bus responds with \activelow{BACK}, \activelow{BERR}, or \activelow{BPER}. If no response is asserted, the CPU holds the
current execution phase and repeats the same bus request on the following clock.

\section{Bus Read/Write Output (BRW)}

\begin{itemize}
	\item \textbf{Direction:} Output
	\item \textbf{Signal Type:} Read High / Write Low
\end{itemize}

Only relevant when the \activelow{BAS} (Bus access strobe output) is asserted.

When high, the CPU is requesting a read from the bus. When low, the CPU is requesting a write to the bus.

\section{Bus Acknowledge Input (BACK)}

\begin{itemize}
	\item \textbf{Direction:} Input
	\item \textbf{Signal Type:} Active Low
\end{itemize}

When asserted, the bus is reporting that the requested read/write completed successfully. On reads, the CPU samples
D0--D15 on the clock edge that accepts \activelow{BACK}.

If all bus devices are fast and can complete reads/writes in a single cycle, this pin can be permanently held asserted.

\section{Bus Error Input (BERR)}

\begin{itemize}
	\item \textbf{Direction:} Input
	\item \textbf{Signal Type:} Active Low
\end{itemize}

When asserted, the bus is reporting that the requested read/write was unsuccessful.

This causes the CPU to abort the I/O operation and raise a bus fault exception, except during an exception-vector
fetch. The response to a failed vector fetch is described in the Exception and Interrupt Handling chapter of the
\archmanual{}: a failed fault-vector fetch raises a double fault, and a failed reset-vector fetch re-enters reset.

For example, a bus error might be raised by I/O to an address that is not connected to any device.

\section{Bus Protection Error Input (BPER)}

\begin{itemize}
	\item \textbf{Direction:} Input
	\item \textbf{Signal Type:} Active Low
\end{itemize}

When asserted, the bus is reporting that the requested read/write was to a valid address/device, but that the I/O is
disallowed given the current state of the CPU or bus.

This causes the CPU to abort the I/O operation and raise a bus protection fault exception, except during an
exception-vector fetch. The response to a failed vector fetch is described in the Exception and Interrupt Handling
chapter of the \archmanual{}: a failed fault-vector fetch raises a double fault, and a failed reset-vector fetch
re-enters reset.

It behaves much like the \activelow{BERR} pin but raises a different fault so that programs can handle it differently.

For example, a bus protection error might be raised by a write to ROM, or by a write to a privileged section of memory
when the CPU is not in supervisor mode.

\section{Bus Access Type Bits (BAT0--BAT2)}

\begin{itemize}
	\item \textbf{Direction:} Output
	\item \textbf{Signal Type:} Active High
\end{itemize}

Output by the CPU to describe what I/O operation it is performing.

These pins can provide hints to external devices during I/O.

For example, if a memory controller knows that the DMA coprocessor will be reading more than one sequential address, it
can prefetch the next address.

Another example is a memory controller providing memory protection when these pins are used with the PROT
(Protected-mode output) pin. A memory controller might assert \activelow{BPER} (bus protection error) if BAT is 010
(data read), PROT is 1 (protected mode) and the address is in a protected memory region.

\begin{table}[H]
	\centering
	\begin{tabular}{rrrl}
		\toprule
		\textbf{BAT2} & \textbf{BAT1} & \textbf{BAT0} & \textbf{Description}        \\
		\midrule
		0             & 0             & 0             & No I/O                      \\
		0             & 0             & 1             & Instruction Fetch           \\
		0             & 1             & 0             & Data Read                   \\
		0             & 1             & 1             & Data Write                  \\
		1             & 0             & 0             & Exception Vector Fetch      \\
		1             & 0             & 1             & DMA Coprocessor Read Burst  \\
		1             & 1             & 0             & DMA Coprocessor Write Burst \\
		1             & 1             & 1             & Reserved                    \\
		\bottomrule
	\end{tabular}
	\caption{Bus Access Type Bits Meaning}
\end{table}

\section{Protected-Mode Output (PROT)}

\begin{itemize}
	\item \textbf{Direction:} Output
	\item \textbf{Signal Type:} Active High
\end{itemize}

Indicates whether the CPU is in protected mode (high) or in supervisor mode (low).

Follows the status register Protected Mode (P) bit, except that it is deasserted for every exception-vector fetch (BAT
100), which is a supervisor access even when the exception is taken from protected mode.

When the CPU is in protected mode, programs have access to fewer types of instructions, for example they cannot write
to the high word of the address registers.

Exposing this as an external pin allows other devices such as memory controllers to block access to protected regions
of memory by raising bus protection errors.

\section{External CPU Halt Input (HALT)}

\begin{itemize}
	\item \textbf{Direction:} Input
	\item \textbf{Signal Type:} Active Low
\end{itemize}

When asserted, the CPU stops executing instructions.

\activelow{HALT} is sampled at every phase 0. If it is asserted there, the CPU stops at that phase 0 with the bus idle. The halt
can occur between an instruction and its coprocessor dispatch slot, between DMA slots, or before a fault, interrupt, or
reset dispatch. Halt latency is at most six cycles. The CPU performs no bus cycles while halted. The interrupt lines
are still sampled and a Level 5 edge is still latched, but nothing is dispatched until the CPU resumes. The CPU resumes
exactly where it stopped when the input is deasserted. A reset ends the halt, but if \activelow{HALT} is still asserted after
reset, the CPU halts at the first phase 0, before the reset-vector fetch.

This pin is useful for hardware-driven debugging, as the CPU can be halted and the state of the system examined.

\section{Force Trace Mode Input (TRCE)}

\begin{itemize}
	\item \textbf{Direction:} Input
	\item \textbf{Signal Type:} Active Low
\end{itemize}

When asserted at an instruction boundary, the instruction being started is traced regardless of whether software has
cleared the Trace Mode (T) flag in the status register. No trace fault is raised while a fault handler runs (exception
level 7), since it would escalate to a double fault.

In trace mode, the CPU raises an exception on every instruction to allow software debugging.

\section{Maskable Interrupt Inputs (IRQ1--IRQ4)}

\begin{itemize}
	\item \textbf{Direction:} Input
	\item \textbf{Signal Type:} Active Low
\end{itemize}

When asserted, the corresponding hardware interrupt line is sampled by the CPU.

The lines can be masked by the Hardware Interrupt Enable Bits (E1--E4) of the status register.

Interrupt inputs are level-sensitive. They are sampled on every clock cycle except during reset: from the cycle in
which reset is requested until the reset-vector fetch completes, no interrupt line, including \activelow{NMI}, is
sampled. Enabled interrupt lines are latched in the pending hardware-exceptions register; disabled maskable lines are
ignored and not queued, and a line is not latched while its own handler or a higher-priority one is running. The
requesting device must hold the line asserted until its handler acknowledges it; an assertion withdrawn before it is
latched is not recognized. If an interrupt pin remains asserted after its handler returns, it may be sampled again
through the normal interrupt rules.

See the Exception and Interrupt Handling chapter of the \archmanual{} for more details.

\section{Non-Maskable Interrupt Inputs (NMI)}

\begin{itemize}
	\item \textbf{Direction:} Input
	\item \textbf{Signal Type:} Active Low
\end{itemize}

When asserted, the CPU samples the Level 5 hardware exception line. Unlike \activelow{IRQ1}--\activelow{IRQ4}, which
are level-sensitive, \activelow{NMI} is edge-triggered: it fires once on assertion regardless of how long the pin is
held, and the pin must be released and reasserted to fire again. An \activelow{NMI} held asserted across a reset does
not fire afterward, because reset does not clear the edge detector's last sampled level and so no transition occurs;
the pin must be released and reasserted.

This interrupt cannot be masked. A new assertion of this pin (a fresh edge) while a Level 5 handler is running raises a
Level Five Hardware Exception Conflict fault; holding the pin asserted from the edge that started the handler does not.
A further edge while a Level 5 request is still pending (for example, behind a fault handler) is lost and is reported
by a single conflict fault when Level 5 is next selected.

Because \activelow{NMI} is edge-triggered, holding the pin asserted after its handler returns does not by itself cause
it to fire again. The pin must be released (driven high) for at least one CLKI cycle and asserted (driven low) again;
the request fires on the clock edge at which the assertion is first sampled. Level 5 exception handlers must therefore
execute quickly, and \activelow{NMI} should be asserted only for important, non-interruptible events.

\section{Reset Input (RSTI)}

\begin{itemize}
	\item \textbf{Direction:} Input
	\item \textbf{Signal Type:} Active Low
\end{itemize}

External devices can use this pin to force a reset of the CPU.

When asserted, the CPU immediately aborts the current instruction or bus wait, stops normal processing, resets to phase
0, and begins the reset sequence described in the Exception and Interrupt Handling chapter of the \archmanual{}.

The \activelow{RSTO} pin is also asserted.

The reset-output hold is six CLKI cycles. The cycle in which \activelow{RSTI} is first sampled asserted counts as the
first, and \activelow{RSTO} stays asserted for five further cycles after \activelow{RSTI} is sampled deasserted, so an
\activelow{RSTI} assertion of $n$ cycles gives $n + 5$ cycles of \activelow{RSTO}. A software reset and a reset-vector
fetch retry each give a six-cycle hold.

When this pin is deasserted, the CPU completes the reset-output hold, stops asserting the \activelow{RSTO} pin, and the
exception unit fetches the reset vector.

If a system needs external devices to be reset for longer than the CPU reset-output hold, reset timing must be managed
externally to the CPU, and \activelow{RSTI} must remain asserted for as long as the other devices need to reset.

\section{Reset Output (RSTO)}

\begin{itemize}
	\item \textbf{Direction:} Output
	\item \textbf{Signal Type:} Active Low
\end{itemize}

Asserted while the \activelow{RSTI} pin is asserted and during the reset-output hold before reset-vector fetch.

It is also asserted for the reset-output hold after a software reset, after a reset caused by a failed double-fault
vector fetch, and for a further hold before the CPU retries a reset-vector fetch that failed. Each of these holds is
six CLKI cycles. A reset caused by a failed double-fault vector fetch behaves as a software reset: the six-cycle
\activelow{RSTO} hold is followed by the reset-vector fetch.

A reset can be triggered by asserting the Reset input (\activelow{RSTI}) pin, or in software by the exception
coprocessor meta-instruction \mnemonic{RSET}.

This pin is useful to reset other devices on the bus whenever the CPU is reset or to trigger more complex reset/boot
controllers.
